problem:  
Let \((M^n,g)\) be a closed oriented Riemannian manifold, and let \(u:M\to S^{m-1}\) be a smooth map.  
Consider energies of the form  
\[
E(u,g) = \int_M \big(|du|_g^p + \alpha\,R_g\,|du|_g^{p-2} + \beta\,\langle \operatorname{Ric}_g(\nabla u,\nabla u)\rangle_g\big)\, d\mathrm{vol}_g,
\]
where \(R_g\) and \(\operatorname{Ric}_g\) are scalar and Ricci curvatures, and \(\alpha,\beta\) are constants.  
It is known that for \(p=n\), the first term is conformally invariant in itself.  The second and third terms, however, typically destroy conformal invariance because curvature transforms nontrivially under \(g\mapsto e^{2\psi}g\).

**Problem:**  
1. Classify all *local geometric Lagrangians* of the broad form  
   \[
   L(u,g)=\text{(polynomial in $du$ and curvature tensors up to order $\le2$)},
   \]
   such that the corresponding Euler–Lagrange equation  
   \(\tau_L(u;g)=0\) is conformally covariant in the sense
   \[
   \tau_L(u; e^{2\psi}g) = e^{-w\psi} \,\tau_L(u;g)
   \]
   for some conformal weight \(w\).  
2. Determine whether any such Lagrangians exist *beyond* the well‑known conformally invariant examples (e.g. the \(n\)-harmonic equation in all dimensions and the conformally invariant wave‑map equation in \(n=2\)).  
3. If no further examples exist, provide a rigorous geometric obstruction explaining why local invariants built from \(du\) and curvature cannot achieve conformal covariance—e.g., via a uniqueness theorem analogous to the Graham–Jenne–Mason–Sparling (GJMS) classification for scalar conformally covariant operators.

Why is it a "good" problem:  
This reformulation lifts the discussion from verifying energy invariance to classifying **all conformally covariant elliptic systems depending locally on curvature and first derivatives of maps into spheres**.  It matches the spirit of known classification problems in conformal geometry (like GJMS operators) but in a genuinely *nonlinear, vector‑valued* setting where no theory currently exists.  Determining even uniqueness would unify ideas from geometric analysis, invariant theory, and nonlinear PDEs.  
It is simple to state—seeking all possible local metrics–map couplings that remain conformally covariant—yet its depth lies in relating transformation laws, tensorial calculus, and elliptic regularity; a complete solution would open a new chapter in the interaction between **conformal geometry and nonlinear elliptic systems**, satisfying all hallmarks of a “good” research‑level problem.